\chapter{Results Chapter Template}
\label{ch:results-template}

% ======================================================================================
% The mini-thesis scaffold for a results chapter (thesis-writer.md, "Results-chapter
% scaffold"). Copy this file to <NN>_<slug>.tex, keep the anchors, replace the text.
%
%   MANDATORY anchors (reviewer raises [MAJOR] when missing):
%     per experiment:   \paragraph{Hypothesis.}
%     chapter closing:  \paragraph{Limitations.}  \paragraph{Bridge.}
%   RECOMMENDED:        \paragraph{Lens.}  \paragraph{Factors.}
%                       \paragraph{Domain conclusion.}  \paragraph{Method conclusion.}
%   OPTIONAL:           \paragraph{Chapter baseline.}
%
% Every number in this chapter traces to an evidence note in docs/evidence_notes/ with
% Status: run-complete or interpreted. A pending value is written as \todo{value pending}.
% ======================================================================================

% --- Chapter intro (no \section header): 1 to 3 paragraphs -----------------------------
% Which factors this chapter studies; why they are interesting (domain and/or method
% lens); a brief roadmap; optionally a one-sentence preview of the answer.

\todo{Chapter intro: factors under study, why they matter, roadmap.} The chapter opens
by naming the factors it studies and the question it answers. It then previews the
structure, so that the reader knows which experiment answers which part.

\section{Experimental setup}
\label{sec:ch2-setup}

% Shared baseline (factors not under study held at their defaults); a design table per
% experiment; the evaluation metric; the seed policy; one line per experiment to come.

\todo{Shared baseline, design table, metric, seed policy.} The baseline configuration
holds every factor not under study at its default value. The metric and the seed policy
are stated once here and not repeated in the experiment sections.

\section{Experiment 1}
\label{sec:ch2-experiment-1}

% Opening paragraph (no heading): motivation for THIS experiment, distinct from the
% chapter intro.

\todo{Motivation for this experiment, from the domain or the method.} The paragraph
states what this experiment adds beyond the chapter intro. It ends by naming the factor
that is swept.

\subsection{Setup}
\label{sec:ch2-experiment-1-setup}

\paragraph{Lens.} \todo{Performance, interpretability, efficiency or robustness: one
sentence per lens.}

\paragraph{Factors.} \todo{What is swept (identifier, options) and what is held fixed.}

\paragraph{Hypothesis.} (i)~\todo{A falsifiable claim, with its reason as a clause.}
(ii)~\todo{A second claim, if the experiment tests one.}

\subsection{Results}
\label{sec:ch2-experiment-1-results}

% Facts only. One paragraph per lens, in Setup order. Every figure and table named via
% \ref before it appears. No interpretation here.

\todo{Results per lens, referencing figures and tables.} The paragraph reports the
numbers, the trends and the patterns that the figures show. It does not interpret them.

\subsection{Discussion}
\label{sec:ch2-experiment-1-discussion}

% Per hypothesis, in the same numbered order as Setup: confirmed, partly confirmed or
% refuted against the seed spread, then one or two sentences on the "why".

\todo{Verdict per hypothesis against the seed spread, then the soft why.} Hypothesis (i)
is addressed first and hypothesis (ii) second. A result within the seed spread is
reported as indistinguishable from the baseline at this seed budget.

\section{Discussion}
\label{sec:ch2-discussion}

\todo{Two or three sentences synthesising the experiments.} The synthesis states the
headline finding of the chapter. It connects the experiments to the question posed in the
chapter intro.

\paragraph{Domain conclusion.} \todo{What the findings mean in the domain, hedged.}

\paragraph{Method conclusion.} \todo{What the findings mean for the method; quantify the
headline effect once, here.}

\paragraph{Limitations.} \todo{Three to five concrete caveats: seed budget, controls not
run, scope of the data or the model, factors not separable in this design.}

\paragraph{Bridge.} \todo{What is carried forward to the next chapter, and what is
deferred.}
